\section{A benchmark with known sources}
\label{sec:benchmark}

\paragraph{Generator.}
We train a shallow piecewise-linear recurrent neural network (shPLRNN; \citealp{brenner2024shPLRNN,hess2023gtf}) on the 22-channel BCI Competition IV-2a EEG \citep{brunner2008bci,tangermann2012review}, with $M$ latent states and a learned linear observation model $x = Wz$.
The trained model is used strictly as a generator, in the sense of a formal surrogate of the system it was fitted to \citep{durstewitz2023reconstructing}, the use that recurrent latent-dynamics models of neural recordings are built for \citep{pandarinath2018lfads,koppe2019identifying,ramezanian2022generative}: its stochastic free run is the ground-truth source signal $s$ and its $W$ the mixing.
One generator is trained per source count $M \in \{2, 4, 8, 16, 22, 32\}$; the clean operating point is $M = 4$ sources in $D = 22$ channels.
Training on EEG supplies the mixing geometry (a learned $22 \times M$ observation matrix whose conditioning grows with $M$), a latent dimensionality per $M$ and process noise calibrated on the recordings; the stochastic free run is spectrally flatter than the latents the model encodes from EEG, and Appendix~\ref{app:realism} reruns the benchmark with a recipe whose free run matches the EEG spectrum.
Every cell draws a training and a held-out trajectory of $T = 2000$ samples (8~s at 250~Hz, cut into 20 segments of 0.4~s) from the same generator; methods are fit on the first and scored on the second (Appendix~\ref{app:protocol}).

\paragraph{Stresses.}
Each axis violates one assumption in a controlled way, harder to the right in Figure~\ref{fig:axes}: (a) a $1/f$ background, the aperiodic component of every EEG spectrum \citep{donoghue2020parameterizing}, at decreasing signal-to-noise ratio; (b) a pointwise nonlinearity after mixing; (c) first-order coupling between the sources; (d) the diversity of the segment-wise modulation, from each source following its own variance profile to one gain shared by all; (e) the number of sources against the fixed channel count, past the point where sources outnumber channels.
Axis (d) is parametrised by an observable quantity, the \textit{co-modulation fraction}: the fraction of the segment-to-segment change of the channel covariance that follows a single spatial pattern (Appendix~\ref{app:placement}).
It can be computed on any recording, which is what lets real EEG be placed on the axis in Section~\ref{sec:real}.

\paragraph{Methods.}
PCA and FastICA \citep{hyvarinen1999fastica} are the linear baselines; FastICA is the positive control for linear mixing of independent sources.
TCL and iVAE identify through segment-wise nonstationarity, PCL through temporal dependence, and CEBRA \citep{schneider2023cebra} is a current neural-data method whose time-contrastive objective likewise rests on local temporal structure.
All methods receive the same 20 segments and return $M$ components; the iVAE is trained on per-channel standardised inputs with the evidence lower bound under a unit-variance Gaussian likelihood.
Implementation details are given in Appendix~\ref{app:methods}.

\paragraph{Scores and references.}
Methods that return source estimates are scored by the mean correlation with the signed sources after each component is matched to its best source (MCC; \citealp{hyvarinen2016tcl,khemakhem2020ivae}).
TCL and CEBRA return energy-like features (features that track the power of a source rather than its signed waveform), since under variance modulation TCL identifies $s^2$ up to a linear map (Appendix~\ref{app:tcl}); they are scored by the canonical-correlation MCC against $|s|$ (which is unchanged by any linear map of the features), and the two scores are read within a row of Figure~\ref{fig:axes}, never across.
Two references accompany every score.
The \textit{chance band} is the 95th percentile of the score of surrogate components that are statistically matched to the sources but independent of them, measured per cell.
The \textit{rotation reference} is the score of the true sources under a random rotation, i.e.\ of a method that recovers the source subspace but does not separate it; a permutation MCC at that level is subspace recovery, not identification (Appendix~\ref{app:scores}).
Error bars throughout are bootstrap 95\% intervals over seeds.
